\scanpage{67}
\margin{\uline{Cas de} $\mathcal O_{X_{\mathrm{strat}}}$\par\uline{et} $\mathcal O_{X_{\mathrm{cris}}}$}
6.4. Cas $F=\mathcal O_{X_{\mathrm{strat}}}$ ou $\mathcal O_{X_{\mathrm{cris}}}$. Alors $\mathcal F^*$ est un \uline{Anneau cosimplicial}\inserted{sur $X_{\mathrm{zar}}$}. \deleted{\unsure{Dans l'étude suivante}}
Dans le cas de 6.2. c'est la limite projective des anneaux cosimpliciaux $\mathcal F_{(i)}^*$, où $\mathcal F_{(i)}^*$ est celui attaché à $Y(i)$ (NB. $Y(i)_{\mathrm{zar}}=X_{\mathrm{zar}}$ !). On est donc, pour l'étude de cet \unsure{hypercoh}, ramené essentiellement au \deleted{contexte} \unsure{stratifiant}.

Chaque $\mathcal F^{(\nu)}$\inserted{sur $X_{\mathrm{zar}}$} est une $\mathcal O_{X_{\mathrm{zar}}}$-algèbre en $\nu+1$ façons canoniques, grâce aux $\nu+1$ projections $\mathrm{pr}_i$ ($\unsure{1\le i\le\nu+1}$) de $X^{\nu+1}$ sur $X$ donc $X^{\nu+1}_{/X}$ sur $X$. Si on choisit p.ex. $\mathrm{pr}_1$, \inserted{$\mathcal F^{(\nu)}$ devient pro-quasi-cohérent.}\deleted{les} opérations simpliciales, et par suite ces opérateurs cobords doivent être considérés comme \uline{op. diff. d'ordre infini} : si on tronque les $\mathcal F^{(\unsure{\nu})}$, on trouvera des op. diff ordinaires ..... \uline{C'est via ces complexes de cochaines canoniques} (plutôt que par les complexes de De Rham). \uline{que les complexes}
\mspage{66}
